\documentclass[a4paper,11pt]{article}

\newcommand{\TPnumero}{Project 4}
\newcommand{\macrosPath}{../../preamble/macros.tex}
% =========================================================
% Tutorial / lab preamble — Time Series Analysis module
% article class + fancyhdr + tcolorbox + listings (English)
% Author : Aymen Ben Brik (aymen.benbrik@esprit.tn)
% =========================================================

\usepackage[utf8]{inputenc}
\usepackage[T1]{fontenc}
\usepackage[english]{babel}
\usepackage{lmodern}
\usepackage{amsmath, amssymb, amsthm, mathtools, bm}
\usepackage{graphicx}
\usepackage{booktabs}
\usepackage{array}
\usepackage{multirow}
\usepackage{colortbl}
\usepackage{xcolor}
\usepackage{tikz}
\usetikzlibrary{positioning, arrows.meta, calc, shapes, fit, backgrounds, matrix}
\usepackage{listings}
\usepackage[most]{tcolorbox}
\usepackage{geometry}
\geometry{a4paper, margin=2cm, headheight=15pt}
\usepackage{fancyhdr}
\usepackage[hidelinks]{hyperref}
\usepackage{enumitem}

% --- Colours ---
\definecolor{espritBlue}{RGB}{30, 60, 110}
\definecolor{espritAccent}{RGB}{200, 80, 30}
\definecolor{boxEx}{RGB}{30, 60, 110}
\definecolor{boxQ}{RGB}{180, 120, 0}
\definecolor{boxC}{RGB}{30, 100, 60}
\definecolor{boxAtt}{RGB}{200, 80, 30}

% --- Listings ---
\definecolor{codebg}{RGB}{248,248,248}
\definecolor{codeframe}{RGB}{220,220,220}
\definecolor{keyword}{RGB}{0,82,155}
\definecolor{comment}{RGB}{88,110,117}
\definecolor{string}{RGB}{133,153,0}

\lstdefinestyle{pythonstyle}{
  language=Python,
  basicstyle=\ttfamily\small,
  keywordstyle=\color{keyword}\bfseries,
  commentstyle=\color{comment}\itshape,
  stringstyle=\color{string},
  backgroundcolor=\color{codebg},
  frame=single,
  rulecolor=\color{codeframe},
  frameround=tttt,
  breaklines=true,
  showstringspaces=false,
  tabsize=2
}
\lstdefinestyle{Rstyle}{
  language=R,
  basicstyle=\ttfamily\small,
  keywordstyle=\color{keyword}\bfseries,
  commentstyle=\color{comment}\itshape,
  stringstyle=\color{string},
  backgroundcolor=\color{codebg},
  frame=single,
  rulecolor=\color{codeframe},
  frameround=tttt,
  breaklines=true,
  showstringspaces=false,
  tabsize=2
}
\lstset{style=pythonstyle}

% --- Common macros (configurable path) ---
\providecommand{\macrosPath}{../../preamble/macros.tex}
\input{\macrosPath}

% --- Exercise counter ---
\newcounter{excounter}
\renewcommand{\theexcounter}{\arabic{excounter}}

\newtcolorbox[use counter=excounter]{exercise}[1][]{
  enhanced, breakable,
  colback=boxEx!5, colframe=boxEx, fonttitle=\bfseries,
  title={Exercise~\theexcounter\ifstrempty{#1}{}{ -- #1}},
  arc=2pt, boxrule=0.6pt
}
\newtcolorbox{question}[1][]{
  enhanced, breakable,
  colback=boxQ!5, colframe=boxQ, fonttitle=\bfseries,
  title={Question\ifstrempty{#1}{}{ -- #1}},
  arc=2pt, boxrule=0.5pt
}
\newtcolorbox{solution}[1][]{
  enhanced, breakable,
  colback=boxC!5, colframe=boxC, fonttitle=\bfseries,
  title={Solution\ifstrempty{#1}{}{ -- #1}},
  arc=2pt, boxrule=0.5pt
}
\newtcolorbox{warning}[1][]{
  enhanced, breakable,
  colback=boxAtt!5, colframe=boxAtt, fonttitle=\bfseries,
  title={Warning\ifstrempty{#1}{}{ -- #1}},
  arc=2pt, boxrule=0.5pt
}

% --- Headers / footers ---
\pagestyle{fancy}
\fancyhf{}
\fancyhead[L]{\textsc{\moduleNom} -- \auteurNom}
\fancyhead[R]{\TPnumero}
\fancyfoot[C]{\thepage}
\fancyfoot[L]{\auteurInstitution}
\fancyfoot[R]{\auteurEmail}
\renewcommand{\headrulewidth}{0.4pt}
\renewcommand{\footrulewidth}{0.2pt}

% --- Placeholder ---
\providecommand{\TPnumero}{Lab}


\title{Project 4: Full Box--Jenkins study with unit-root testing\\
\large ADF/KPSS $+$ differencing $+$ SARIMA on real monthly data}
\author{\auteurNom\\\auteurEmail\\\auteurInstitution}
\date{\anneeUniversitaire}

\begin{document}
\maketitle

\section*{Context}

Project 3 left us with a non-white-noise residual after a deterministic
log-decomposition of \texttt{SouvenirSales.csv}. In this project you
push the analysis to its conclusion: you classify the series as TS or
DS using formal unit-root tests, apply the appropriate transformation
(detrending or differencing), fit a SARIMA model, and produce a
twelve-month forecast.

The deliverable is a complete, reproducible Box--Jenkins study that
will serve as a template for Chapter 5 (ARMA) and Chapter 6
(volatility).

\section*{Deliverables}
\begin{itemize}
  \item Reproducible \texttt{project4\_<name>.py} with seed $42$.
  \item Report PDF (6--8 pages) including all figures, all numerical
        outputs, and a final forecast plot.
  \item Optional 8-min oral defence.
\end{itemize}

% =========================================================
\begin{exercise}[Visual diagnosis and transformation]
\textbf{(1)} Load \texttt{data/SouvenirSales.csv}. Plot the raw
series $Y_t$. Comment on (i) the trend, (ii) the multiplicative
seasonality (amplitude growing with the level).

\textbf{(2)} Take logs: $y_t = \log Y_t$. Plot $y_t$. Verify
visually that the seasonality is now \emph{additive} (constant
amplitude).

\textbf{(3)} On $y_t$, plot the sample ACF up to lag $36$. Identify
the seasonal lags ($12, 24, 36$) and the slow decay typical of a
non-stationary series.
\end{exercise}

% =========================================================
\begin{exercise}[Unit-root testing]
\textbf{(1)} Run the ADF test on $y_t$ with three deterministic
specifications:
\begin{itemize}
  \item \texttt{regression="n"}: no constant, no trend;
  \item \texttt{regression="c"}: constant only;
  \item \texttt{regression="ct"}: constant and trend.
\end{itemize}
For each, report stat, $p$-value, $5\%$ critical value, decision.

\textbf{(2)} Run the KPSS test on $y_t$ with \texttt{regression="c"}
and \texttt{regression="ct"}. Report stat, $p$-value, decision.

\textbf{(3)} Build the ADF/KPSS combination table. Which cell does
$y_t$ fall in? (TS, DS, or ambiguous?)
\end{exercise}

% =========================================================
\begin{exercise}[Differencing]
\textbf{(1)} Compute the regular difference $\Delta y_t = y_t -
y_{t-1}$. Plot $\Delta y_t$ and its ACF up to lag $36$.

\textbf{(2)} The seasonal lags should still dominate. Apply the
seasonal difference $\Delta_{12} = (1 - L^{12})$ on top:
\[
  w_t = \Delta_{12} \Delta y_t = (1 - L)(1 - L^{12}) y_t.
\]
Plot $w_t$ and its ACF/PACF up to lag $36$.

\textbf{(3)} Run ADF and KPSS on $w_t$ (regression \texttt{c}).
Confirm: $w_t$ is stationary --- $y_t$ is $\text{SARIMA}(\,\cdot\,,1,\,\cdot\,)
(\,\cdot\,,1,\,\cdot\,)_{12}$.
\end{exercise}

% =========================================================
\begin{exercise}[SARIMA model selection]
\textbf{(1)} From the ACF/PACF of $w_t$, propose two candidate
SARIMA orders. The classical airline-style choice is
$(0,1,1)(0,1,1)_{12}$ but you may also consider, e.g.,
$(1,1,0)(1,1,0)_{12}$ or $(1,1,1)(0,1,1)_{12}$.

\textbf{(2)} Fit each candidate by maximum likelihood. Report:
\begin{itemize}
  \item estimated coefficients (with standard errors);
  \item AIC and BIC;
  \item Ljung--Box $p$-value at lag $24$ on residuals;
  \item residual standard deviation.
\end{itemize}

\textbf{(3)} Pick the best model on the (AIC, BIC, Ljung--Box) triplet.
Justify your choice in 5 lines.
\end{exercise}

% =========================================================
\begin{exercise}[Forecasting and back-transformation]
\textbf{(1)} On the chosen model, produce a 12-month forecast of
$\log Y_t$ with $95\%$ confidence intervals.

\textbf{(2)} Back-transform: $\widehat Y_{T+h} = \exp(\widehat
y_{T+h} + \widehat\sigma^2/2)$ if you want the conditional mean of
$Y$, or simply $\exp(\widehat y_{T+h})$ for the conditional median.
Plot the original $Y_t$ together with the back-transformed forecast
and CI on the original scale.

\textbf{(3)} Discuss the interpretation: what does the model predict
for sales in December of the year following 2001?
\end{exercise}

% =========================================================
\begin{exercise}[Bonus: hold-out validation]
Split the data into a training set (1995--2000, 72 observations) and
a test set (2001, 12 observations).
\begin{itemize}
  \item Fit the chosen SARIMA on the training set.
  \item Compute the forecast errors on the test set.
  \item Report RMSE and MAPE on the original scale.
\end{itemize}
Discuss: does the model generalise out of sample?
\end{exercise}

% =========================================================
\section*{Marking grid (out of 100)}

\begin{center}
\begin{tabular}{p{8cm}rl}
\toprule
\textbf{Criterion} & \textbf{Points} & \\
\midrule
Visual diagnosis + log transform              & 10 \\
ADF and KPSS tests + combination matrix       & 20 \\
Differencing $(1-L)(1-L^{12})$ + verification & 15 \\
SARIMA candidates + AIC/BIC/LB selection      & 25 \\
Forecast + back-transform                     & 15 \\
Bonus: hold-out validation                    & 10 \\
Reproducibility + report quality              &  5 \\
\midrule
\textbf{Total}                                & \textbf{100} \\
\bottomrule
\end{tabular}
\end{center}

\bigskip
\hfill\auteurNom\par
\hfill\auteurEmail\par

\end{document}
